\subsection{Wurzeln als Potenzen}
\label{subsec:fundamentals_roots_as_powers}

Die Schreibweise für rationale Exponenten soll mit dem dritten Potenzgesetz
verträglich sein. Für $a>0$ setzen wir daher $\sqrt[n]{a}=a^r$. Da die
$n$-te Potenz der $n$-ten Wurzel wieder $a$ ergibt, muss gelten:
\begin{ClozeArea}{3}
  \centering
  $\displaystyle
    (a^r)^n=a
    \quad\Longleftrightarrow\quad
    a^{rn}=a^1
    \quad\Longleftrightarrow\quad
    rn=1
    \quad\Longleftrightarrow\quad
    r=\frac1n.$
\end{ClozeArea}

\begin{KeyPointBox}{Rationale Exponenten}{fundamentals_rational_exponents}
  Für $a>0$, $m\in\Z$ und $n\in\N$ mit $n\geq2$ gilt
  \[
    \sqrt[n]{a}=a^{1/n},
    \qquad
    a^{m/n}=\sqrt[n]{a^m}=\left(\sqrt[n]{a}\right)^m.
  \]
\end{KeyPointBox}

\begin{KeyPointBox}{Wurzelgesetze}{fundamentals_root_laws}
  Da sich Wurzeln als Potenzen schreiben lassen, gelten analog zu den Potenzgesetzen (kurz: PG) die folgenden Wurzelgesetze.
  Sie sind keine neuen Rechengesetze, sondern unmittelbare Folge der Potenzgesetze.
  \begin{align*}
    \sqrt[n]{a}\,\sqrt[n]{b} =\sqrt[n]{ab},
    &\qquad \text{, da} \quad
    \Cloze{12cm}{\sqrt[n]{a}\,\sqrt[n]{b} = a^{1/n}b^{1/n} \overset{\text{2. PG}}{=} (ab)^{1/n} =\sqrt[n]{ab}}; \\[0.9em]
    %
    \frac{\sqrt[n]{a}}{\sqrt[n]{b}} =\sqrt[n]{\frac ab},
    &\qquad \text{, da} \quad
    \Cloze{12cm}{\frac{\sqrt[n]{a}}{\sqrt[n]{b}} =\frac{a^{1/n}}{b^{1/n}} \overset{\text{2. PG}}{=} \left(\frac ab\right)^{1/n} =\sqrt[n]{\frac ab}}; \\[0.9em]
    %
    \sqrt[n]{\sqrt[m]{a}} =\sqrt[nm]{a},
    &\qquad \text{, da} \quad
    \Cloze{12cm}{\sqrt[n]{\sqrt[m]{a}} =\left(a^{1/m}\right)^{1/n} \overset{\text{3. PG}}{=} a^{1/(mn)} =\sqrt[nm]{a}}.
  \end{align*}
\end{KeyPointBox}

\begin{RemarkBox}[remark title=Wurzeln bei Produkten und Summen]
  Wurzeln dürfen über Produkte und Quotienten verteilt werden, aber keinesfalls
  über Summen oder Differenzen.

  \begin{center}
    \begin{tabularx}{0.94\linewidth}{@{}>{\centering\arraybackslash}X >{\centering\arraybackslash}X@{}}
      $\sqrt{ab}\,\Cloze{8mm}{=}\,\sqrt a\cdot\sqrt b$
        \IfClozeTF{}{%
          \quad\textcolor{green!50!black}{%
            \tikz[baseline=-0.55ex]
              \draw[line width=1.3pt]
                (0,0.6ex) -- (0.5ex,0) -- (1.55ex,1.4ex);%
          }%
        }
      &
      \textcolor{red!75!black}{%
        $\sqrt{a+b}\,\Cloze{8mm}{\neq}\,\sqrt a+\sqrt b$%
      }
        \IfClozeTF{}{%
          \quad\textcolor{red!75!black}{\textbf{i.\,A. falsch}}%
        }
    \end{tabularx}
  \end{center}
\end{RemarkBox}

\begin{ExampleBox}{Rationale Exponenten}{fundamentals_rational_exponents_example}
  \begin{enumerate}[label=\alph*),itemsep=0.35em]
    \item $27^{2/3} = \Cloze{12cm}{\left(\sqrt[3]{27}\right)^2=3^2=9}$
    \item $9^{-1/2} = \Cloze{12cm}{\displaystyle\frac1{\sqrt9}=\frac13}$
    \item $4^{2,5} = \Cloze{12cm}{4^{\frac{5}{2}}=\left(4^{\frac{1}{2}}\right)^5=\left(\sqrt{4}\right)^5=4\sqrt{2}}$
    \item $\sqrt{\sqrt[3]{64}} = \Cloze{12cm}{\sqrt[6]{64}=2}$
  \end{enumerate}
\end{ExampleBox}
